%% annex-b-type-catalog.tex -- generated by tools/build_from_docx.py; edit freely

\hypertarget{annex-b-initial-content-of-the-type-catalog}{%
\section{Initial Content of the Type Catalog}\label{annex-b-initial-content-of-the-type-catalog}}

The type catalog of Section 3.6 is a registry of versioned entries to
which descriptors refer by type-ref, in the form identifier@version.
This annex fixes the fields of each kind of entry and the content with
which the catalog is initialized. The reason codes and verification
methods, which T2.3 owns, are given in full. The source kinds,
controller bindings, metrics, control parameters and action types are
initial entries of the domains that own them, and the objective KPIs
those of T2.1; each owner extends its part of the catalog by registering
further entries under the same structure.

\hypertarget{b.1-entry-structure}{%
\subsection{Entry Structure}\label{b.1-entry-structure}}

Every entry carries an identifier, a version under Semantic Versioning
2.0.0, its entry kind, its owner, its conformance class --- normative,
or extension where the identifier carries the prefix
\texttt{x-\textless{}publisher\textgreater{}-} --- and an informative
description. An entry version referenced by an admitted descriptor is
immutable. Table B-1 states the fields each entry kind carries in
addition.

\begin{wptable}{B-1}

\begin{longtable}[]{@{}
  >{\raggedright\arraybackslash}p{(\columnwidth - 2\tabcolsep) * \real{0.2125}}
  >{\raggedright\arraybackslash}p{(\columnwidth - 2\tabcolsep) * \real{0.7875}}@{}}
\caption{Fields of the catalog entries by kind}\tabularnewline
\toprule\noalign{}
\begin{minipage}[b]{\linewidth}\raggedright
\textbf{Entry kind}
\end{minipage} & \begin{minipage}[b]{\linewidth}\raggedright
\textbf{Fields in addition to identifier, version, kind, owner,
conformance and description}
\end{minipage} \\
\midrule\noalign{}
\endfirsthead
\toprule\noalign{}
\begin{minipage}[b]{\linewidth}\raggedright
\textbf{Entry kind}
\end{minipage} & \begin{minipage}[b]{\linewidth}\raggedright
\textbf{Fields in addition to identifier, version, kind, owner,
conformance and description}
\end{minipage} \\
\midrule\noalign{}
\endhead
\bottomrule\noalign{}
\endlastfoot
\textbf{Metric} & unit; measurement\_\allowbreak{}kind (Table 4-1);
object\_\allowbreak{}class; definition; formula, for a computed metric;
source\_\allowbreak{}kinds; aggregations, the permitted subset of those
of the kind. \\
\textbf{Control parameter} & object\_\allowbreak{}class;
value\_\allowbreak{}type and unit; controller\_\allowbreak{}bindings
through which it is set; readback\_\allowbreak{}method.\allowbreak{} \\
\textbf{Action type} & object\_\allowbreak{}class;
control\_\allowbreak{}parameters changed;
parameter\_\allowbreak{}schema; default\_\allowbreak{}idempotency;
verification\_\allowbreak{}methods available;
rollback\_\allowbreak{}method; controller\_\allowbreak{}mappings, the
controller operation for each controller binding;
unverifiable\_\allowbreak{}rate\_\allowbreak{}ceiling.\allowbreak{} \\
\textbf{Insight type} & schema; semantics;
subject\_\allowbreak{}object\_\allowbreak{}class.\allowbreak{} \\
\textbf{Objective KPI} & definition; unit; methodology, in the form used
by D2.1; project\_\allowbreak{}kpis.\allowbreak{} \\
\textbf{Verification method} & readback\_\allowbreak{}path,
ms\_\allowbreak{}observation or
controller\_\allowbreak{}state\_\allowbreak{}query;
typical\_\allowbreak{}latency.\allowbreak{} \\
\textbf{Controller binding} & controller\_\allowbreak{}kind; interface
and interface\_\allowbreak{}version; action\_\allowbreak{}types
supported; adapter, the southbound adapter of the DMO. \\
\textbf{Source kind} & protocol; read\_\allowbreak{}only, always
true. \\
\textbf{Reason code} & category: decision, mediation, actuation or
lifecycle; meaning. \\
\textbf{Service contract} & contribution\_\allowbreak{}class; provider;
interface and interface\_\allowbreak{}version;
consumption\_\allowbreak{}point, for a platform service. \\
\end{longtable}

\end{wptable}

\hypertarget{b.2-measurement-kinds-and-permitted-aggregations}{%
\subsection{Measurement Kinds and Permitted Aggregations}\label{b.2-measurement-kinds-and-permitted-aggregations}}

The measurement kind of a metric entry fixes which aggregations a
descriptor may declare for it (Table 4-1). Table B-2 states the
permitted values of \texttt{functions.MS.metrics{[}{]}.aggregation} for
each kind. Admission rejects any other value, and the runtime applies
the treatment of the kind in computing the aggregate.

\begin{wptable}{B-2}

\begin{longtable}[]{@{}
  >{\raggedright\arraybackslash}p{(\columnwidth - 4\tabcolsep) * \real{0.2250}}
  >{\raggedright\arraybackslash}p{(\columnwidth - 4\tabcolsep) * \real{0.3125}}
  >{\raggedright\arraybackslash}p{(\columnwidth - 4\tabcolsep) * \real{0.4625}}@{}}
\caption{Permitted aggregations by measurement kind}\tabularnewline
\toprule\noalign{}
\begin{minipage}[b]{\linewidth}\raggedright
\textbf{Measurement kind}
\end{minipage} & \begin{minipage}[b]{\linewidth}\raggedright
\textbf{Permitted aggregations}
\end{minipage} & \begin{minipage}[b]{\linewidth}\raggedright
\textbf{Treatment by the runtime}
\end{minipage} \\
\midrule\noalign{}
\endfirsthead
\toprule\noalign{}
\begin{minipage}[b]{\linewidth}\raggedright
\textbf{Measurement kind}
\end{minipage} & \begin{minipage}[b]{\linewidth}\raggedright
\textbf{Permitted aggregations}
\end{minipage} & \begin{minipage}[b]{\linewidth}\raggedright
\textbf{Treatment by the runtime}
\end{minipage} \\
\midrule\noalign{}
\endhead
\bottomrule\noalign{}
\endlastfoot
\textbf{gauge} & none, min, max, mean, p50, p95, p99 & Summation over
time is not permitted. \\
\textbf{cumulative\_\allowbreak{}counter} & none, sum & sum is the sum
of increments; a decrease is a reset, and the window containing it is
reported indeterminate. \\
\textbf{delta\_\allowbreak{}counter} & none, sum & Increments are
summed; no reset handling is required. \\
\textbf{event} & none, count & Occurrences are counted by timestamp. \\
\textbf{state} & none, count, duration & count gives the number of
transitions; duration gives the time spent in each state value within
the window. \\
\textbf{log} & none, count & Structured fields extracted from records
are published as metrics of their own kinds. \\
\textbf{computed} & As permitted for the kind of the result & Evaluated
from the formula of the entry or from the declared derivation (Section
4.4). \\
\end{longtable}

\end{wptable}

\hypertarget{b.3-source-kinds}{%
\subsection{Source Kinds}\label{b.3-source-kinds}}

A source kind names a protocol through which a Monitoring System
acquires data, and guarantees that the acquisition is read-only: the
runtime subscribes, polls or reads, and holds no write path through it.
Table B-3 lists the initial source kinds.

\begin{wptable}{B-3}

\begin{longtable}[]{@{}
  >{\raggedright\arraybackslash}p{(\columnwidth - 4\tabcolsep) * \real{0.3125}}
  >{\raggedright\arraybackslash}p{(\columnwidth - 4\tabcolsep) * \real{0.3750}}
  >{\raggedright\arraybackslash}p{(\columnwidth - 4\tabcolsep) * \real{0.3126}}@{}}
\caption{Initial source kinds}\tabularnewline
\toprule\noalign{}
\begin{minipage}[b]{\linewidth}\raggedright
\textbf{Identifier}
\end{minipage} & \begin{minipage}[b]{\linewidth}\raggedright
\textbf{Protocol}
\end{minipage} & \begin{minipage}[b]{\linewidth}\raggedright
\textbf{Typical use}
\end{minipage} \\
\midrule\noalign{}
\endfirsthead
\toprule\noalign{}
\begin{minipage}[b]{\linewidth}\raggedright
\textbf{Identifier}
\end{minipage} & \begin{minipage}[b]{\linewidth}\raggedright
\textbf{Protocol}
\end{minipage} & \begin{minipage}[b]{\linewidth}\raggedright
\textbf{Typical use}
\end{minipage} \\
\midrule\noalign{}
\endhead
\bottomrule\noalign{}
\endlastfoot
\textbf{source.\allowbreak{}e2\_\allowbreak{}kpm} & O-RAN E2 service
model KPM, by subscription through the Near-RT RIC & Performance
measurements of the radio access network \\
\textbf{source.\allowbreak{}controller\_\allowbreak{}telemetry} &
Telemetry interface of the domain controller, such as that of
TeraFlowSDN & Transport state, counters and alarms \\
\textbf{source.gnmi} & gNMI subscription & Streaming telemetry of
network elements \\
\textbf{source.snmp} & SNMP polling and notifications & Counters and
state of network elements \\
\textbf{source.\allowbreak{}prometheus} & Prometheus scrape endpoint &
Metrics of cloud-native workloads and platform services \\
\textbf{source.\allowbreak{}ngsi\_\allowbreak{}ld} & NGSI-LD context
broker subscription & Context information and data products of the Data
Continuum \\
\end{longtable}

\end{wptable}

\hypertarget{b.4-verification-methods}{%
\subsection{Verification Methods}\label{b.4-verification-methods}}

Two verification methods are defined. Unverifiable is a declaration of
the action space rather than a method: an action type declared
unverifiable executes only within the rate ceiling for unverifiable
execution of its catalog entry, and its outcomes are reported as
acknowledged\_\allowbreak{}unverified (Chapter 7). Table B-4 states the
two methods.

\begin{wptable}{B-4}

\begin{longtable}[]{@{}
  >{\raggedright\arraybackslash}p{(\columnwidth - 4\tabcolsep) * \real{0.2750}}
  >{\raggedright\arraybackslash}p{(\columnwidth - 4\tabcolsep) * \real{0.4750}}
  >{\raggedright\arraybackslash}p{(\columnwidth - 4\tabcolsep) * \real{0.2501}}@{}}
\caption{Verification methods}\tabularnewline
\toprule\noalign{}
\begin{minipage}[b]{\linewidth}\raggedright
\textbf{Identifier}
\end{minipage} & \begin{minipage}[b]{\linewidth}\raggedright
\textbf{Read-back path}
\end{minipage} & \begin{minipage}[b]{\linewidth}\raggedright
\textbf{Typical latency}
\end{minipage} \\
\midrule\noalign{}
\endfirsthead
\toprule\noalign{}
\begin{minipage}[b]{\linewidth}\raggedright
\textbf{Identifier}
\end{minipage} & \begin{minipage}[b]{\linewidth}\raggedright
\textbf{Read-back path}
\end{minipage} & \begin{minipage}[b]{\linewidth}\raggedright
\textbf{Typical latency}
\end{minipage} \\
\midrule\noalign{}
\endhead
\bottomrule\noalign{}
\endlastfoot
\textbf{ms\_\allowbreak{}observation} & The effect is confirmed in the
Monitoring System evidence for the target, independently of the
controller. & One MS window after the change settles \\
\textbf{controller\_\allowbreak{}state\_\allowbreak{}query} & The ACT
runtime reads the configured state of the target from the controller and
compares it with the intended state. & One controller round trip \\
\end{longtable}

\end{wptable}

\hypertarget{b.5-controller-bindings}{%
\subsection{Controller Bindings}\label{b.5-controller-bindings}}

A controller binding names the controller through which an Actuator
executes and the interface it uses; the southbound adapter that realizes
it is supplied by the DMO of the domain. Table B-5 lists the initial
bindings.

\begin{wptable}{B-5}

\begin{longtable}[]{@{}
  >{\raggedright\arraybackslash}p{(\columnwidth - 4\tabcolsep) * \real{0.2875}}
  >{\raggedright\arraybackslash}p{(\columnwidth - 4\tabcolsep) * \real{0.4625}}
  >{\raggedright\arraybackslash}p{(\columnwidth - 4\tabcolsep) * \real{0.2499}}@{}}
\caption{Initial controller bindings}\tabularnewline
\toprule\noalign{}
\begin{minipage}[b]{\linewidth}\raggedright
\textbf{Identifier}
\end{minipage} & \begin{minipage}[b]{\linewidth}\raggedright
\textbf{Controller and interface}
\end{minipage} & \begin{minipage}[b]{\linewidth}\raggedright
\textbf{Domain}
\end{minipage} \\
\midrule\noalign{}
\endfirsthead
\toprule\noalign{}
\begin{minipage}[b]{\linewidth}\raggedright
\textbf{Identifier}
\end{minipage} & \begin{minipage}[b]{\linewidth}\raggedright
\textbf{Controller and interface}
\end{minipage} & \begin{minipage}[b]{\linewidth}\raggedright
\textbf{Domain}
\end{minipage} \\
\midrule\noalign{}
\endhead
\bottomrule\noalign{}
\endlastfoot
\textbf{binding.\allowbreak{}tfs\_\allowbreak{}nbi} & TeraFlowSDN
northbound interface & transport \\
\textbf{binding.\allowbreak{}e2sm\_\allowbreak{}rc} & E2SM-RC through
the Near-RT RIC & oran \\
\textbf{binding.o1} & O1 configuration management of O-RAN managed
elements & oran \\
\textbf{binding.\allowbreak{}dmo\_\allowbreak{}lcm} & Lifecycle
interface of the DMO, over OSM or the domain\textquotesingle s own
orchestration & every domain \\
\textbf{binding.\allowbreak{}idmo\_\allowbreak{}delegation} & Delegation
interface of the IDMO; the only binding of a global agent (V-25) &
idmo \\
\end{longtable}

\end{wptable}

\hypertarget{b.6-reason-codes}{%
\subsection{Reason Codes}\label{b.6-reason-codes}}

A reason code is placed in every rationale, verdict, execution envelope
and lifecycle record, and names why the decision, the verdict, the
outcome or the transition occurred. The set is closed within each
category: a contributor that needs a code the catalog lacks declares an
extension code, which makes its descriptor of extension conformance.
Table B-6 states the initial codes.

\begin{wptable}{B-6}

\begin{longtable}[]{@{}
  >{\raggedright\arraybackslash}p{(\columnwidth - 4\tabcolsep) * \real{0.4000}}
  >{\raggedright\arraybackslash}p{(\columnwidth - 4\tabcolsep) * \real{0.1375}}
  >{\raggedright\arraybackslash}p{(\columnwidth - 4\tabcolsep) * \real{0.4625}}@{}}
\caption{Reason codes}\tabularnewline
\toprule\noalign{}
\begin{minipage}[b]{\linewidth}\raggedright
\textbf{Code}
\end{minipage} & \begin{minipage}[b]{\linewidth}\raggedright
\textbf{Category}
\end{minipage} & \begin{minipage}[b]{\linewidth}\raggedright
\textbf{Meaning}
\end{minipage} \\
\midrule\noalign{}
\endfirsthead
\toprule\noalign{}
\begin{minipage}[b]{\linewidth}\raggedright
\textbf{Code}
\end{minipage} & \begin{minipage}[b]{\linewidth}\raggedright
\textbf{Category}
\end{minipage} & \begin{minipage}[b]{\linewidth}\raggedright
\textbf{Meaning}
\end{minipage} \\
\midrule\noalign{}
\endhead
\bottomrule\noalign{}
\endlastfoot
\textbf{decision.\allowbreak{}objective\_\allowbreak{}at\_\allowbreak{}risk}
& decision & The status of the objective is at-risk, and the intent
seeks to restore it. \\
\textbf{decision.\allowbreak{}objective\_\allowbreak{}breached} &
decision & The objective is not met. \\
\textbf{decision.\allowbreak{}optimization} & decision & The objective
is met, and the intent improves it within bounds. \\
\textbf{decision.\allowbreak{}safe\_\allowbreak{}intent} & decision &
The declared safe intent was emitted under the failure policy. \\
\textbf{decision.\allowbreak{}redecision} & decision & The intent
re-decides after a rejected, unverified or unresolved outcome. \\
\textbf{decision.\allowbreak{}objective\_\allowbreak{}unachievable} &
decision & The re-decision budget is exhausted; a lifecycle event is
raised instead of an intent. \\
\textbf{mediation.\allowbreak{}no\_\allowbreak{}contention} & mediation
& allow: the intent touches only an action space held exclusively. \\
\textbf{mediation.\allowbreak{}within\_\allowbreak{}grant} & mediation &
allow: the intent lies within a valid action envelope or distributed
policy. \\
\textbf{mediation.\allowbreak{}bounded\_\allowbreak{}alternative} &
mediation & modify: a bounded alternative, or a permitted candidate, is
returned. \\
\textbf{mediation.\allowbreak{}contention} & mediation & defer:
contending intents are being resolved. \\
\textbf{mediation.\allowbreak{}priority} & mediation & deny: an intent
of higher effective priority holds the resource. \\
\textbf{mediation.\allowbreak{}policy} & mediation & deny: a governance
policy forbids the action. \\
\textbf{mediation.\allowbreak{}stale} & mediation & deny: the validity
window of the intent has expired. \\
\textbf{mediation.\allowbreak{}retired\_\allowbreak{}generation} &
mediation & deny: the submitting generation is retired. \\
\textbf{mediation.\allowbreak{}scope\_\allowbreak{}violation} &
mediation & The submission exceeds the registered control scope;
recorded as a conformance fault rather than returned as a verdict. \\
\textbf{actuation.\allowbreak{}verified} & actuation &
acknowledged\_\allowbreak{}verified: the change is confirmed by
read-back. \\
\textbf{actuation.\allowbreak{}translation\_\allowbreak{}out\_\allowbreak{}of\_\allowbreak{}bounds}
& actuation & rejected: the intent cannot be translated within the
registered bounds. \\
\textbf{actuation.\allowbreak{}precondition\_\allowbreak{}failed} &
actuation & rejected: a declared precondition does not hold on the
target. \\
\textbf{actuation.\allowbreak{}controller\_\allowbreak{}rejected} &
actuation & rejected: the controller refused the operation. \\
\textbf{actuation.\allowbreak{}token\_\allowbreak{}invalid} & actuation
& rejected: the conflict-resolution token is missing, invalid or issued
for another intent. \\
\textbf{actuation.\allowbreak{}grant\_\allowbreak{}void} & actuation &
rejected: the last-mile check found the grant void. \\
\textbf{actuation.\allowbreak{}rate\_\allowbreak{}limited} & actuation &
rejected: a rate limit, cooldown or minimum hold time applies. \\
\textbf{actuation.\allowbreak{}readback\_\allowbreak{}mismatch} &
actuation & acknowledged\_\allowbreak{}unverified: the read-back state
differs from the intended state. \\
\textbf{actuation.\allowbreak{}readback\_\allowbreak{}unavailable} &
actuation & acknowledged\_\allowbreak{}unverified: no read-back was
obtained within the verification deadline. \\
\textbf{actuation.\allowbreak{}timeout} & actuation & unresolved: no
acknowledgement within the actuation timeout. \\
\textbf{lifecycle.\allowbreak{}gate\_\allowbreak{}failed} & lifecycle &
A check of the activation gate failed; the instance remains
admitted\_\allowbreak{}not\_\allowbreak{}active.\allowbreak{} \\
\textbf{lifecycle.\allowbreak{}demoted} & lifecycle & The authority
state was lowered on a declared threshold. \\
\textbf{lifecycle.\allowbreak{}restored} & lifecycle & Authority was
restored by a decided act with a record. \\
\textbf{lifecycle.\allowbreak{}binding\_\allowbreak{}fault} & lifecycle
& suspended: a binding failed or left its accepted version range. \\
\textbf{lifecycle.\allowbreak{}mediation\_\allowbreak{}loss} & lifecycle
& suspended: the mediation binding was lost, at the expiry of any action
envelope then held. \\
\textbf{lifecycle.\allowbreak{}generation\_\allowbreak{}retired} &
lifecycle & The generation was retired at a handover or an update. \\
\textbf{lifecycle.\allowbreak{}terminated} & lifecycle & The instance
was terminated. \\
\end{longtable}

\end{wptable}

\hypertarget{b.7-metrics}{%
\subsection{Metrics}\label{b.7-metrics}}

Table B-7 lists the initial metric entries of the radio, transport and
service domains. A computed metric is evaluated by the runtime from its
formula, written in the operators of Section 4.4 with the aggregation of
each operand named inside it. The end-to-end key performance indicators
of 3GPP TS 28.554 enter the catalog as computed metrics of the same
form, each citing the clause of that specification from which its
formula is taken.

\begin{wptable}{B-7}

\begin{longtable}[]{@{}
  >{\raggedright\arraybackslash}p{(\columnwidth - 6\tabcolsep) * \real{0.3625}}
  >{\raggedright\arraybackslash}p{(\columnwidth - 6\tabcolsep) * \real{0.1625}}
  >{\raggedright\arraybackslash}p{(\columnwidth - 6\tabcolsep) * \real{0.1000}}
  >{\raggedright\arraybackslash}p{(\columnwidth - 6\tabcolsep) * \real{0.3750}}@{}}
\caption{Initial metrics}\tabularnewline
\toprule\noalign{}
\begin{minipage}[b]{\linewidth}\raggedright
\textbf{Identifier}
\end{minipage} & \begin{minipage}[b]{\linewidth}\raggedright
\textbf{Kind}
\end{minipage} & \begin{minipage}[b]{\linewidth}\raggedright
\textbf{Unit}
\end{minipage} & \begin{minipage}[b]{\linewidth}\raggedright
\textbf{Definition or formula}
\end{minipage} \\
\midrule\noalign{}
\endfirsthead
\toprule\noalign{}
\begin{minipage}[b]{\linewidth}\raggedright
\textbf{Identifier}
\end{minipage} & \begin{minipage}[b]{\linewidth}\raggedright
\textbf{Kind}
\end{minipage} & \begin{minipage}[b]{\linewidth}\raggedright
\textbf{Unit}
\end{minipage} & \begin{minipage}[b]{\linewidth}\raggedright
\textbf{Definition or formula}
\end{minipage} \\
\midrule\noalign{}
\endhead
\bottomrule\noalign{}
\endlastfoot
\textbf{metric.\allowbreak{}oran.\allowbreak{}prb\_\allowbreak{}used\_\allowbreak{}dl,
metric.\allowbreak{}oran.\allowbreak{}prb\_\allowbreak{}used\_\allowbreak{}ul}
& gauge & count & Mean number of physical resource blocks used in the
reporting period, per cell, as measured under 3GPP TS 28.552 and
reported through E2 KPM. \\
\textbf{metric.\allowbreak{}oran.\allowbreak{}prb\_\allowbreak{}available\_\allowbreak{}dl,
metric.\allowbreak{}oran.\allowbreak{}prb\_\allowbreak{}available\_\allowbreak{}ul}
& gauge & count & Mean number of physical resource blocks available in
the reporting period, per cell. \\
\textbf{metric.\allowbreak{}oran.\allowbreak{}prb\_\allowbreak{}utilization\_\allowbreak{}dl,
metric.\allowbreak{}oran.\allowbreak{}prb\_\allowbreak{}utilization\_\allowbreak{}ul}
& computed & ratio & ratio(mean(prb\_\allowbreak{}used),
mean(prb\_\allowbreak{}available)) over the window, per cell and
direction. \\
\textbf{metric.\allowbreak{}oran.\allowbreak{}admission\_\allowbreak{}requests}
& delta\_\allowbreak{}counter & count & Requests for radio resources
received, such as QoS flow setup attempts. \\
\textbf{metric.\allowbreak{}oran.\allowbreak{}admission\_\allowbreak{}refusals}
& delta\_\allowbreak{}counter & count & Requests for radio resources
refused for lack of resources. \\
\textbf{metric.\allowbreak{}oran.\allowbreak{}blocking\_\allowbreak{}probability}
& computed & ratio & ratio(sum(admission\_\allowbreak{}refusals),
sum(admission\_\allowbreak{}requests)) over the window. \\
\textbf{metric.\allowbreak{}transport.\allowbreak{}rsl} & gauge & dBm &
Received signal level at a radio link end. \\
\textbf{metric.\allowbreak{}transport.\allowbreak{}link\_\allowbreak{}capacity}
& gauge & Mbit/s & Capacity of the link at its active modulation. \\
\textbf{metric.\allowbreak{}transport.\allowbreak{}modulation} & state &
--- & Active modulation profile of the link. \\
\textbf{metric.\allowbreak{}transport.\allowbreak{}tx\_\allowbreak{}octets}
& cumulative\_\allowbreak{}counter & byte & Octets transmitted on the
link interface. \\
\textbf{metric.\allowbreak{}transport.\allowbreak{}link\_\allowbreak{}throughput}
& computed & Mbit/s & scale(rate(tx\_\allowbreak{}octets), 8e-6). \\
\textbf{metric.\allowbreak{}transport.\allowbreak{}link\_\allowbreak{}utilization}
& computed & ratio & ratio(link\_\allowbreak{}throughput,
link\_\allowbreak{}capacity).\allowbreak{} \\
\textbf{metric.\allowbreak{}transport.\allowbreak{}one\_\allowbreak{}way\_\allowbreak{}delay}
& gauge & ms & One-way delay of probe packets over a link or a path, by
active probe. \\
\textbf{metric.\allowbreak{}service.\allowbreak{}delivery\_\allowbreak{}latency}
& gauge & ms & End-to-end delivery latency of the packets of a
service. \\
\textbf{metric.\allowbreak{}service.\allowbreak{}jitter} & computed & ms
& difference(p99(delivery\_\allowbreak{}latency),
p50(delivery\_\allowbreak{}latency)) over the window. \\
\textbf{metric.\allowbreak{}service.\allowbreak{}deadline\_\allowbreak{}misses}
& delta\_\allowbreak{}counter & count & Packets of a service delivered
after their deadline. \\
\textbf{metric.\allowbreak{}service.\allowbreak{}state} & state & --- &
Operational state of a service: up, degraded or down. \\
\textbf{metric.\allowbreak{}service.\allowbreak{}availability} &
computed & ratio & Fraction of the window spent in state up, from the
duration aggregation of
metric.\allowbreak{}service.\allowbreak{}state.\allowbreak{} \\
\end{longtable}

\end{wptable}

\hypertarget{b.8-objective-kpis}{%
\subsection{Objective KPIs}\label{b.8-objective-kpis}}

Objective KPIs are the indicators against which objectives are declared
(Section 3.5.6). Table B-8 lists the initial entries and the project
KPIs to which each is linked; their measurement methodology is that of
D2.1.

\begin{wptable}{B-8}

\begin{longtable}[]{@{}
  >{\raggedright\arraybackslash}p{(\columnwidth - 6\tabcolsep) * \real{0.3125}}
  >{\raggedright\arraybackslash}p{(\columnwidth - 6\tabcolsep) * \real{0.0875}}
  >{\raggedright\arraybackslash}p{(\columnwidth - 6\tabcolsep) * \real{0.4375}}
  >{\raggedright\arraybackslash}p{(\columnwidth - 6\tabcolsep) * \real{0.1625}}@{}}
\caption{Initial objective KPIs}\tabularnewline
\toprule\noalign{}
\begin{minipage}[b]{\linewidth}\raggedright
\textbf{Identifier}
\end{minipage} & \begin{minipage}[b]{\linewidth}\raggedright
\textbf{Unit}
\end{minipage} & \begin{minipage}[b]{\linewidth}\raggedright
\textbf{Definition}
\end{minipage} & \begin{minipage}[b]{\linewidth}\raggedright
\textbf{Project KPIs}
\end{minipage} \\
\midrule\noalign{}
\endfirsthead
\toprule\noalign{}
\begin{minipage}[b]{\linewidth}\raggedright
\textbf{Identifier}
\end{minipage} & \begin{minipage}[b]{\linewidth}\raggedright
\textbf{Unit}
\end{minipage} & \begin{minipage}[b]{\linewidth}\raggedright
\textbf{Definition}
\end{minipage} & \begin{minipage}[b]{\linewidth}\raggedright
\textbf{Project KPIs}
\end{minipage} \\
\midrule\noalign{}
\endhead
\bottomrule\noalign{}
\endlastfoot
\textbf{kpi.\allowbreak{}reconfiguration\_\allowbreak{}time} & ms & Time
from the decision to the reconfiguration taking effect in the target
domain. & KPI\#4, KPI\#5 \\
\textbf{kpi.\allowbreak{}control\_\allowbreak{}loop\_\allowbreak{}time}
& ms & Duration of one closed-loop cycle, from observation to verified
actuation. & KPI\#34 \\
\textbf{kpi.\allowbreak{}e2e\_\allowbreak{}latency} & ms & Cross-domain
end-to-end latency. & KPI\#7 \\
\textbf{kpi.jitter} & ms & Cross-domain jitter, with the share of
packets missing their deadline. & KPI\#8 \\
\textbf{kpi.\allowbreak{}service\_\allowbreak{}recovery\_\allowbreak{}time}
& s & Time from the onset of a service-affecting condition to the
recovery of the service. & KPI\#31 \\
\textbf{kpi.\allowbreak{}service\_\allowbreak{}availability} & ratio &
Availability and reliability of the service. & KPI\#36 \\
\textbf{kpi.\allowbreak{}service\_\allowbreak{}continuity} & ratio &
Continuity of the service under AI-driven proactive management. &
KPI\#37 \\
\textbf{kpi.\allowbreak{}conflict\_\allowbreak{}penalty} & ratio &
Performance penalty attributable to conflicts between agents. &
KPI\#21 \\
\textbf{kpi.\allowbreak{}local\_\allowbreak{}resolution\_\allowbreak{}ratio}
& ratio & Share of service management and lifecycle tasks resolved by
local AE and DE. & KPI\#26 \\
\textbf{kpi.\allowbreak{}ai\_\allowbreak{}deployment\_\allowbreak{}time}
& s & Deployment time of an AI-based resource or energy management
algorithm. & KPI\#41 \\
\textbf{kpi.\allowbreak{}bus\_\allowbreak{}delivery\_\allowbreak{}latency}
& ms & End-to-end message delivery latency on the service bus. &
KPI\#24 \\
\textbf{kpi.\allowbreak{}communication\_\allowbreak{}overhead} & ratio &
Management communication overhead relative to a baseline, at maintained
SLA. & KPI\#23 \\
\end{longtable}

\end{wptable}

\hypertarget{b.9-control-parameters-and-action-types-of-the-transport-domain}{%
\subsection{Control Parameters and Action Types of the Transport Domain}\label{b.9-control-parameters-and-action-types-of-the-transport-domain}}

The initial control parameters and action types of the wireless
transport domain are given here as the pattern that the other domains
follow when they register theirs. Each is set through the TeraFlowSDN
northbound interface, and none of the action types admits unverifiable
execution. Table B-9 states them.

\begin{wptable}{B-9}

\begin{longtable}[]{@{}
  >{\raggedright\arraybackslash}p{(\columnwidth - 6\tabcolsep) * \real{0.3499}}
  >{\raggedright\arraybackslash}p{(\columnwidth - 6\tabcolsep) * \real{0.1375}}
  >{\raggedright\arraybackslash}p{(\columnwidth - 6\tabcolsep) * \real{0.2875}}
  >{\raggedright\arraybackslash}p{(\columnwidth - 6\tabcolsep) * \real{0.2250}}@{}}
\caption{Initial control parameters and action types of the transport
domain}\tabularnewline
\toprule\noalign{}
\begin{minipage}[b]{\linewidth}\raggedright
\textbf{Identifier}
\end{minipage} & \begin{minipage}[b]{\linewidth}\raggedright
\textbf{Entry kind}
\end{minipage} & \begin{minipage}[b]{\linewidth}\raggedright
\textbf{Content}
\end{minipage} & \begin{minipage}[b]{\linewidth}\raggedright
\textbf{Verification and rollback}
\end{minipage} \\
\midrule\noalign{}
\endfirsthead
\toprule\noalign{}
\begin{minipage}[b]{\linewidth}\raggedright
\textbf{Identifier}
\end{minipage} & \begin{minipage}[b]{\linewidth}\raggedright
\textbf{Entry kind}
\end{minipage} & \begin{minipage}[b]{\linewidth}\raggedright
\textbf{Content}
\end{minipage} & \begin{minipage}[b]{\linewidth}\raggedright
\textbf{Verification and rollback}
\end{minipage} \\
\midrule\noalign{}
\endhead
\bottomrule\noalign{}
\endlastfoot
\textbf{param.\allowbreak{}transport.\allowbreak{}acm\_\allowbreak{}range}
& Control parameter & Minimum and maximum modulation profiles between
which adaptive coding and modulation operates on a link. & Read back by
controller\_\allowbreak{}state\_\allowbreak{}query.\allowbreak{} \\
\textbf{param.\allowbreak{}transport.\allowbreak{}tx\_\allowbreak{}power}
& Control parameter & Transmit power of a radio link end, in dBm. & Read
back by
controller\_\allowbreak{}state\_\allowbreak{}query.\allowbreak{} \\
\textbf{param.\allowbreak{}transport.\allowbreak{}path\_\allowbreak{}route}
& Control parameter & Ordered links of an explicit path. & Read back by
controller\_\allowbreak{}state\_\allowbreak{}query.\allowbreak{} \\
\textbf{action.\allowbreak{}transport.\allowbreak{}set\_\allowbreak{}acm\_\allowbreak{}range}
& Action type & Changes
param.\allowbreak{}transport.\allowbreak{}acm\_\allowbreak{}range on a
link; idempotent; mapped to a configuration update of the link
endpoints. & controller\_\allowbreak{}state\_\allowbreak{}query or
ms\_\allowbreak{}observation; rollback restores the previous range. \\
\textbf{action.\allowbreak{}transport.\allowbreak{}set\_\allowbreak{}tx\_\allowbreak{}power}
& Action type & Changes
param.\allowbreak{}transport.\allowbreak{}tx\_\allowbreak{}power on a
link end; idempotent. &
controller\_\allowbreak{}state\_\allowbreak{}query or
ms\_\allowbreak{}observation; rollback restores the previous power. \\
\textbf{action.\allowbreak{}transport.\allowbreak{}reroute\_\allowbreak{}path}
& Action type & Replaces
param.\allowbreak{}transport.\allowbreak{}path\_\allowbreak{}route of a
path; idempotent. & controller\_\allowbreak{}state\_\allowbreak{}query;
rollback restores the previous route. \\
\end{longtable}

\end{wptable}
